\section{My team}

\textbf{2a. Structure --- discussed once, drifted twice, re-opened under pressure.}
At the wk12 kickoff we did hold an explicit structure conversation, and I want to credit Robin for forcing it: I would have skipped it. We split on module ownership --- Robin took the flight loop and state execution, I took perception and simulation, Edward took the Cube and hardware interface, Demetro took client-facing communication and the FDR presentation, and our PM coordinated schedule and budget. We deliberately declined to name a single technical lead because the brief (R08) devolves autonomy-level choices to the team and we wanted that authority distributed. The structure then drifted twice. The first drift was silent: around wk16, as hardware kept slipping, the project's centre of gravity migrated from the bench to my laptop because that was the only surface where forward progress was possible --- nobody declared it, it just happened. The second drift was loud: at wk17, after I caught Robin and myself building overlapping state machines in parallel, I asked for a renegotiation, and we agreed in writing (\texttt{docs/DESIGN\_DECISIONS.md:107--109}) that his script would own the flight loop and call my \texttt{passive\_watch} as an HTTP service for detections. Looking back I now see we optimised for task division over communication cadence because task division felt unambiguous and cadence felt like overhead --- and the first silent drift happened precisely because we had no routine for restating structure when the environment shifted. I did not understand before this project that unexplicit structure is a luxury which only survives in slack-rich teams.

\textbf{2b. What went well --- and the wk17 FSM conflict I had to manage.}
The most impactful team-working element, measured by what it unblocked, was a handoff ratio I had not previously noticed: the four hours I spent writing \texttt{docs/VISION\_PIPELINE\_FOR\_COLLEAGUE.md} and the \texttt{--robin} integration flag (commit \texttt{a5b1ab7}) bought Robin roughly twenty hours of unblocked integration time across wk20 --- he wired \texttt{passive\_watch} into his mission script without asking me a single clarifying question for two days. I now understand those docs worked because they traded my solo time for Robin's onboarding time at roughly 5:1, better team economics than I had assumed documentation ever delivered; I had previously treated handoff docs as a cost I paid for politeness. The second element was the Demetro FDR collaboration: I rewrote his two speaking scripts four times (commits \texttt{b2a9cd4, 9b826f8, 7920e5a, 794ef5c}) until his delivery rehearsal stopped drifting off the 3-minute target, and he told me afterwards he finally felt confident presenting the autonomy argument.

The conflict I want to name happened in wk17. I had proposed an 11-state FSM for \texttt{main.py} with explicit transitions, timeout guards, and a geofence repulsor; Robin and Demetro pushed back in the same Teams thread, Robin arguing that a two-out-of-five maintainability floor meant the FSM was over-engineered, Demetro arguing the extra states would not survive a hardware-debug session with Edward. My first reaction was defensive --- I genuinely believed the FSM was safer given our NFZ edge cases, and commit \texttt{ce5c036} later proved at least part of that belief correct when the RTL mode tuple fix needed an explicit state transition to \texttt{LANDING}. But I took 24 hours, reread Robin's flight script, and realised half of my insistence was that the FSM was my code. I proposed a compromise on the wk17 call: I would keep the full FSM for the autonomous path in \texttt{main.py}, but strip it to a lightweight polling structure exposed over HTTP for Robin's script --- that is what the \texttt{--robin} flag in \texttt{passive\_watch\_2.py} actually is, not a feature but an interface concession. Robin accepted, Demetro signed off, and the wk20 integration took an afternoon instead of a week. The Hatton-level move I now see is that I had been conflating ``is this the right design'' with ``is this the right design \emph{for this team},'' and these two kinds of rightness can come apart. I was right about the FSM and wrong about the context. The leadership-without-title moment, such as it was, came the following week: when the wk18 flight was cancelled for weather and the group call went quiet, I was the one who said out loud that we should stop waiting for hardware and commit to a simulation-first plan, delegated the SITL home-location fix to Edward, asked Demetro to move his FDR slot forward, and wrote the pivot up in a 400-word Teams message the same night. Nobody had asked me to decide, and nobody objected when I did.

\textbf{2c. Even better if --- three concrete changes, with falsifiable signals.}
First, I would institute a weekly one-line decision log from wk1, not wk12 --- every call that bound more than one person's work, with date, owner, and the reason it was made. The silent wk16 drift would have been caught by any written record; the signal that it was working would be: zero overlapping-work surprises at integration checkpoints. Second, I would front-load interface contracts \emph{before} parallel work begins, specifically a two-hour whiteboard session in the first week of any new module to agree on the inputs, outputs, and failure modes at every seam --- the Robin-vs-me state-machine overlap was exactly the failure mode high-agency individuals default into under stress, and a single wk14 session on the \texttt{passive\_watch} HTTP contract would have removed the whole wk17 conflict. Third, I would stop treating documentation as a substitute for conversation. The two pieces of writing that actually changed teammates' behaviour in this project were the 400-word pivot message and the 2000-word complaint letter --- both worked because they were short, timely, and accompanied by a live call. The 3000-word colleague handoff doc, by contrast, sat unread until Robin finally needed it. The falsifiable signal for next time: any document longer than 500 words must be paired with a 15-minute walkthrough inside a week of being written, or I will treat it as unsent.
